\section{Coloring rounds with dependent entry weights}\label{sec:rounds}

We now use Lemma~\ref{lem:adjustment} to extend a partial correspondence
coloring while controlling the weights of the entries that remain
available.  Entropy and edge energy, together with compensated weight
changes, are features of Johansson's coloring method; see
\cite[Sections~2 and~6--8]{BansalGuptaGuruganesh2018}.  Here the adjustment
operations act on the graph of list entries.  Entries in a single list need not evolve
independently; the bound on the number of entries from each part in an
operation will supply the concentration needed for a simultaneous update.

Fix $r\ge4$ and let $B=B(r)$ be the constant in
Lemma~\ref{lem:adjustment}.  Choose
\begin{equation}\label{eq:round-constants}
 K=16(B+2),\qquad \gamma=0.003,\qquad C=\frac{100K}{\gamma},\qquad
 k=\left\lceil\frac{C\Delta}{\log\Delta}\right\rceil.
\end{equation}
We prove Theorem~\ref{thm:main} with $C_r=C$.  Restrict each list to
exactly $k$ entries, restricting the matchings accordingly, and continue
to write $L(v)$ for the resulting lists.  For sufficiently large
$\Delta$ we have $k\le\Delta$.  Throughout the coloring process, $\Delta$
denotes the maximum degree of the original graph.  Set
\begin{equation}\label{eq:round-parameters}
 w_* = \frac Kk,\qquad \Lambda=20\log\Delta,\qquad
 \tau=\Delta^{-0.10},\qquad \xi=\Delta^{-0.14},\qquad
 \zeta=\Delta^{-0.15}.
\end{equation}
The parameter $\tau$ sets the activation and adjustment rates, $\xi$ is
the permitted error per round, and $\zeta$ is a smaller tolerance used
in the concentration estimates.
All asymptotic estimates below are uniform over the graphs, lists,
matchings, and states satisfying the stated assumptions, with $r$ fixed.
In particular, a factor $\Delta^{o(1)}$ in an upper bound is at most
$\Delta^\varepsilon$ for every fixed $\varepsilon>0$ once $\Delta$ is
sufficiently large in terms of $r$ and $\varepsilon$.  The parameters
$U,P_0$ in Lemma~\ref{lem:adjustment} are $\Delta^{o(1)}$, and
$R_0=O(\log\log\Delta)$, at the value of $\Lambda$ in
\eqref{eq:round-parameters}.

Consider a proper partial correspondence coloring, and let $G'$ be the
subgraph induced by its uncolored vertices.  An entry at an uncolored
vertex is \emph{available} if it conflicts with none of the colors already
chosen at neighboring vertices.  Give each entry $i\in L(v)$, for
$v\in V(G')$, a nonnegative weight $w_i$, requiring every positive-weight
entry to be available.  Let $H$ be the graph on the positive-weight
entries whose edges are the forbidden pairs in the restricted matchings.
Its parts are the sets $L(v)\cap V(H)$.  The map sending an entry to its
host vertex is a graph homomorphism $H\to G'$.  A clique in $H$ has
distinct hosts, so $H$ is $K_r$-free.  The matching condition also implies
that each entry has at most one neighbor in any part.  These are precisely
the structural assumptions of Lemma~\ref{lem:adjustment}.

Use $s_i,d_i,t_i$ for the neighborhood, incident-edge, and triangle
weights in that lemma, calculated in $H$.  We maintain
\begin{equation}\label{eq:weight-bounds}
 \Delta^{-3}\le w_i\le\Delta^{-0.99},\qquad s_i\le\Lambda
 \quad (i\in V(H)).
\end{equation}
For $v\in V(G')$ and $vu\in E(G')$, define
\begin{equation}\label{eq:vertex-potentials}
 \begin{aligned}
 W_v&=\sum_{i\in L(v)}w_i,\\
 D_{vu}&=\sum_{\substack{i\in L(v),\ l\in L(u)\\il\in E(H)}}w_iw_l,
 &D_v&=\sum_{u\in N_{G'}(v)}D_{vu},\\
 F_v&=\sum_{i\in L(v)}w_*f(w_i/w_*)
       +\gamma\log\Delta\,(K-W_v)+D_v.
 \end{aligned}
\end{equation}
The entropy sum includes zero-weight entries, using $f(0)=1$.
Both that sum and $D_v$ are nonnegative, and
$D_v=\sum_{i\in L(v)\cap V(H)}d_i\le\Lambda W_v$.
Thus an upper bound on $F_v$ prevents the total weight $W_v$ from
becoming too small, while coloring neighbors decreases its edge term.

Initially, color no vertices and give every entry weight $w_*$.  Since
$w_*=\Theta(\log\Delta/\Delta)$ and
$s_i\le\Delta w_*\le(K/C)\log\Delta$, the bounds
\eqref{eq:weight-bounds} hold.  Moreover,
\begin{equation}\label{eq:initial-state}
 W_v=K,\qquad F_v=D_v\le \frac{K^2\Delta}{k}
                         \le\frac{K^2}{C}\log\Delta.
\end{equation}
The following proposition is the step we shall iterate.

\begin{proposition}[One coloring round]\label{prop:one-round}
For all sufficiently large $\Delta$ depending only on $r$, suppose that
a proper partial correspondence coloring and entry weights satisfy
availability, \eqref{eq:weight-bounds}, and
\begin{equation}\label{eq:mass-window}
 K/2\le W_v\le2K\qquad(v\in V(G')).
\end{equation}
There is an extension of the partial coloring and a choice of weights
on the remaining entries that preserve availability and
\eqref{eq:weight-bounds}, and satisfy
\begin{equation}\label{eq:round-conclusions}
 W_v^{\mathrm{new}}\le W_v+\xi,\qquad
 F_v^{\mathrm{new}}\le F_v+\xi,\qquad
 \frac{\deg_{\mathrm{new}}(v)}{\Delta}
 \le(1-\tau)\frac{\deg_{G'}(v)}{\Delta}+\xi
\end{equation}
for every vertex $v$ that remains uncolored.
\end{proposition}

\subsection{The random update and its expectation}

Assume the hypotheses of Proposition~\ref{prop:one-round}; the case
$V(G')=\varnothing$ is immediate.  Fix the operations and rates $a_j$
given by Lemma~\ref{lem:adjustment} on $H$, with listed supports $J_j$.
We use two families of mutually independent Bernoulli flags.  For each
entry $z\in V(H)$ an \emph{activation flag} is set with probability
$\tau w_z$; it touches every neighbor of $z$ in $H$, with designated
multiplier zero there.  For each operation $j$ its flag is set with
probability $\tau a_j$ and touches $J_j$.  A multiplier operation has
its prescribed multipliers $m_i^{(j)}$ on $J_j$, and a drop has designated
multiplier zero on its singleton support.  The probabilities are valid
for large $\Delta$: the nonempty support and rate bound of
Lemma~\ref{lem:adjustment} give $a_j\le\Lambda^4$.

For each $i\in V(H)$ put $\bar w_i=w_iY_i$.  If exactly one set flag
touches $i$, let $Y_i$ be its designated multiplier; if at least two set
flags touch $i$, let $Y_i=0$.  If no set flag touches $i$, put
\begin{equation}\label{eq:no-flag-multiplier}
 Y_i=1+\tau\left(s_i+\sum_{j\in\mathcal M}a_j(1-m_i^{(j)})\right),
\end{equation}
where $\mathcal M$ indexes the multiplier operations and their
multipliers are extended by $1$ off the listed supports.  Zero-weight
entries stay at zero.  The coefficient of $\tau$ in
\eqref{eq:no-flag-multiplier} has absolute value at most
$\Lambda+\Lambda^4(1+U)$.  The no-flag multiplier is therefore positive
and is $1+O(\tau\Delta^{o(1)})$.

An entry \emph{activates} when its activation flag is set.  An activating
entry is \emph{successful} if none of its neighbors in $H$ activates.
At every host with a successful entry, choose one such entry as its color,
using any fixed order on that list to break ties.  Successful entries
cannot be adjacent in $H$, and they were available before the round,
so this extends the proper partial coloring.  Every entry conflicting
with a newly chosen color is touched by its activation flag and hence
has provisional weight zero, whether or not other flags also touch it.
Consequently, after newly colored vertices are removed, all remaining
positive provisional weights are available.
The set of hosts receiving a color depends only on activation flags,
independently of the choices among successful entries.

Until that removal, a bar on any quantity in
\eqref{eq:vertex-potentials} means that it is calculated with the
provisional weights on the entire start-of-round graph $G'$.  We may
continue to use $E(H)$ in these sums because zero weights remain zero.
The compensation in \eqref{eq:no-flag-multiplier} cancels the
first-order mean loss from activations and multiplier operations.
The following estimate shows how the adjustment lemma controls the
remaining cost in the potential.

\begin{lemma}[Expected changes]\label{lem:round-expectation}
For the random update just defined,
\begin{equation}\label{eq:round-expectations}
 \begin{aligned}
 \EE\bar W_v&\le W_v+o(\xi),\\
 \EE\bar F_v&\le F_v+\tau(B+1)D_v+o(\xi),\\
 \EE\bar D_{vu}&\ge\tfrac12D_{vu}\qquad(vu\in E(G')).
 \end{aligned}
\end{equation}
\end{lemma}

\begin{proof}
Let $a_i^-$ be the sum of the rates of drops at $i$.  Expanding according
to the set flags touching one or two entries gives the following
identities, with $il\in E(H)$ in the last line:
\begin{equation}\label{eq:multiplier-moments}
 \begin{aligned}
 \EE(Y_i-1)
   &=-\tau a_i^-+O(\tau^2\Delta^{o(1)}),\\
 \EE f(Y_i)
   &=\tau\left(s_i+a_i^-+
           \sum_{j\in\mathcal M}a_jf(m_i^{(j)})\right)
       +O(\tau^2\Delta^{o(1)}),\\
 \EE[(Y_i-1)(Y_l-1)]
   &=\tau\sum_{z\in N_H(i)\cap N_H(l)}w_z\\
   &\quad+\tau\sum_{j\in\mathcal M}
          a_j(m_i^{(j)}-1)(m_l^{(j)}-1)
       +O(\tau^2\Delta^{o(1)}).
 \end{aligned}
\end{equation}
Here is a justification of the errors and the first-order terms.
The sum of the probabilities of flags touching either of two fixed
entries is at most $2\tau(\Lambda+\Lambda^4)$.  Independence bounds
both the probability of two or more such flags and the total error in
replacing the exact single-flag probabilities by their marginal probabilities by
$O(\tau^2\Delta^{o(1)})$.  All multipliers and their $f$-values are
$\Delta^{o(1)}$, so these errors remain of the stated order.
For the first line, the single flags contribute
$-s_i-a_i^-+\sum_{j\in\mathcal M}a_j(m_i^{(j)}-1)$ to the coefficient
of $\tau$; \eqref{eq:no-flag-multiplier} cancels all but $-a_i^-$.
For the second line, use $f(0)=1$ and $f(1+h)=O(h^2)$ at the no-flag
value.  For the third line, only flags touching both entries contribute
to first order.  These are common-neighbor activations and shared
multiplier flags.  A singleton drop cannot touch both endpoints, and
a flag touching only one endpoint has an additional factor
$O(\tau\Delta^{o(1)})$ from the other endpoint.

For every $i\in V(H)$, including when $Y_i=0$, the entropy identity is
\begin{equation}\label{eq:entropy-change}
 w_*\bigl[f(w_iY_i/w_*)-f(w_i/w_*)\bigr]
   =w_i\log(w_i/w_*)(Y_i-1)+w_if(Y_i).
\end{equation}
Allocate the change in $F_v$ to its positive-weight entries.  The
contribution of $i\in L(v)\cap V(H)$ is
\begin{equation}\label{eq:entry-potential-change}
 w_i\bigl[(\log(w_i/w_*)-\gamma\log\Delta)(Y_i-1)+f(Y_i)\bigr]
 +\sum_{l\in N_H(i)}w_iw_l(Y_iY_l-1).
\end{equation}
Its expectation is at most
\begin{equation}\label{eq:entry-expected-change}
 \tau\left(d_i+t_i+\sum_j a_jg_i^{(j)}\right)
 +O\bigl((w_i+d_i)\tau^2\Delta^{o(1)}\bigr).
\end{equation}
Indeed, \eqref{eq:multiplier-moments} gives the activation terms $d_i$
from entropy and $t_i$ from edge correlations.  For a multiplier
operation, the compensated linear terms cancel, leaving exactly
$w_if(m_i^{(j)})$ and the quadratic edge terms that define $g_i^{(j)}$.
A drop at $i$ contributes
$w_i(1-\log(w_i/w_*)+\gamma\log\Delta)\le\Lambda w_i$ to the
non-edge part and $-d_i$ to the edge part.  A drop at a neighbor $l$
contributes $-w_iw_l$.  These are bounded by the drop vector in
Lemma~\ref{lem:adjustment}.  Finally,
$|\log(w_i/w_*)|=O(\log\Delta)$ by \eqref{eq:weight-bounds}; absorbing
these logarithms gives the error in \eqref{eq:entry-expected-change}.

Sum \eqref{eq:entry-expected-change} over $L(v)\cap V(H)$ and use
$t_i+\sum_j a_jg_i^{(j)}\le Bd_i$.  Since
$D_v\le\Lambda W_v$ and $W_v\le2K$, the total error is
$O((W_v+D_v)\tau^2\Delta^{o(1)})=o(\xi)$.  This proves the potential
estimate.  The first line of \eqref{eq:multiplier-moments} proves the
mass estimate.  For each edge $il$ in the sum defining $D_{vu}$, the
probability that neither endpoint is touched by a set flag is
$1-O(\tau\Delta^{o(1)})$; on this event,
$Y_iY_l=1+O(\tau\Delta^{o(1)})$.  Nonnegativity on the other events
therefore gives $\EE(w_iw_lY_iY_l)\ge w_iw_l/2$ for large $\Delta$.
Summing proves the last assertion.
\end{proof}

\subsection{Concentration with bounded input participation}

We next make the expectation bounds simultaneous at all vertices.
The relevant concentration principle is a bounded-variable form of
the generalized H\"older inequality of Finner~\cite{Finner} and the
read-$k$ concentration method of Gavinsky, Lovett, Saks, and
Srinivasan~\cite{GavinskyReadK}.  We give the short proof, including
the exponential-moment estimate of Hoeffding~\cite{Hoeffding1963}.

\begin{lemma}[Bounded input participation]\label{lem:read-concentration}
Let $X_1,\ldots,X_m$ be bounded real functions of finitely many
independent inputs.  Suppose that $X_j$ has range of length at most
$b_j$, and that each input is used by at most $p$ of the functions,
where $p\ge1$ is real.  If $\sum_jb_j^2>0$, then for every $a>0$,
\begin{equation}\label{eq:input-concentration}
 \PP\left(\sum_j(X_j-\EE X_j)\ge a\right),\quad
 \PP\left(\sum_j(X_j-\EE X_j)\le-a\right)
 \ \le\ \exp\left(-\frac{2a^2}{p\sum_jb_j^2}\right).
\end{equation}
If $\sum_jb_j^2=0$, both probabilities are zero.
\end{lemma}

\begin{proof}
For nonnegative bounded functions $Z_j$ with the same dependence
sets, we first prove
\begin{equation}\label{eq:input-holder}
 \EE\prod_j Z_j\le\prod_j(\EE Z_j^p)^{1/p}.
\end{equation}
Induct on the number of independent inputs.  Condition on all inputs
except the last, and let $I$ index the functions using that last input.
H\"older's inequality with exponent $p$ for the $|I|\le p$ factors,
and with a constant factor $1$ of reciprocal exponent $1-|I|/p$ when
$|I|<p$, bounds this conditional integral
by $\prod_j\widetilde Z_j$, where
$\widetilde Z_j=(\EE_{\mathrm{last}}Z_j^p)^{1/p}$ for $j\in I$ and
$\widetilde Z_j=Z_j$ otherwise.  Each $\widetilde Z_j$ uses only the
remaining inputs of $Z_j$, so the participation bound still holds.
The induction hypothesis applies and gives
$\EE\prod_j\widetilde Z_j\le
\prod_j(\EE\widetilde Z_j^p)^{1/p}
=\prod_j(\EE Z_j^p)^{1/p}$, proving \eqref{eq:input-holder}.

If a random variable $X$ has range of length $b$, then
\begin{equation}\label{eq:hoeffding-moment}
 \log\EE e^{s(X-\EE X)}\le\frac{s^2b^2}{8}\qquad(s\in\mathbb R).
\end{equation}
To see this, the second derivative of the logarithm on the left is
the variance under the exponentially tilted law.  Any law supported
on an interval of length $b$ has variance at most $b^2/4$, by measuring
squared distance from the midpoint.  The logarithm and its first
derivative vanish at zero, which proves \eqref{eq:hoeffding-moment}.
Apply \eqref{eq:input-holder} to
$Z_j=e^{s(X_j-\EE X_j)}$ and then \eqref{eq:hoeffding-moment} with
parameter $ps$.  The resulting moment bound is
$\exp(ps^2\sum_jb_j^2/8)$.  Exponential Markov with
$s=4a/(p\sum_jb_j^2)$ proves the upper tail; replacing $X_j$ by
$-X_j$ proves the lower tail.
\end{proof}

Apply Lemma~\ref{lem:read-concentration} first to $\bar W_v$ and to
the non-edge part $\bar F_v-\bar D_v$.  Apart from an additive constant,
each is a sum of at most $k\le\Delta$ functions, one for each entry of
$L(v)$, of range length at most $\Delta^{-0.99+o(1)}$.  This follows
from \eqref{eq:weight-bounds}, \eqref{eq:no-flag-multiplier}, and
$U=\Delta^{o(1)}$; the logarithmic factors in the entropy term are
absorbed in $\Delta^{o(1)}$.  This also covers arbitrarily small
positive provisional weights: with $0\log0=0$,
$\sup_{0\le z\le b}|z\log z|=b\log(1/b)$ for $0<b<e^{-1}$.
A multiplier flag occurs in at most
$P_0$ of these functions, by its support bound in a single part.
A drop flag occurs in at most one, and an activation flag occurs in
at most one because its activating entry has at most one neighbor
in $L(v)$.  Thus $p=P_0$ is valid.  In particular, each requirement
\begin{equation}\label{eq:vertex-concentration}
 \bar W_v\le\EE\bar W_v+\zeta,\qquad
 \bar F_v-\bar D_v\le\EE(\bar F_v-\bar D_v)+\zeta
\end{equation}
fails with probability at most $\exp(-\Delta^{0.5})$ for large $\Delta$:
the exponent in \eqref{eq:input-concentration} is at least
$\Delta^{0.68-o(1)}$.

For an edge $vu\in E(G')$, the matching sum $\bar D_{vu}$ has at
most $k$ terms, each of range length at most
$\Delta^{-1.98+o(1)}$.  Every entry at either endpoint occurs in at
most one term, so each flag occurs in at most $2P_0$ terms.  Applying
Lemma~\ref{lem:read-concentration} with $p=2P_0$ shows that
\begin{equation}\label{eq:matching-concentration}
 |\bar D_{vu}-\EE\bar D_{vu}|\le\Delta^{-1.25}
\end{equation}
fails with probability at most $2\exp(-\Delta^{0.4})$.  Here the
exponent is at least $\Delta^{0.46-o(1)}$.  This separation into
non-edge terms and individual matching sums is what permits the
participation bound to apply; no independence among entries of a
list is assumed.

For ordinary list coloring, operations can instead be chosen separately
in each graph induced by one common color label.  Sampling their flags
independently makes contributions of distinct colors independent, so
the corresponding concentration bound has $p=1$.  Arbitrary edge
matchings provide no such common labels; the per-part support bound
replaces that independence here.

\subsection{Activations that produce colors}

The expectation estimate allows a potential increase of order
$\tau D_v$.  To offset it, we must show that enough neighbors receive
colors, with their contributions measured both by $D_{vu}$ and by
$1/\Delta$.  Fix $v\in V(G')$, and take one of the coefficient families
\begin{equation}\label{eq:neighbor-coefficients}
 b_u=D_{vu}\quad(u\in N_{G'}(v)),\qquad\text{or}\qquad
 b_u=1/\Delta\quad(u\in N_{G'}(v)).
\end{equation}
In either case $0\le b_u\le2K\Delta^{-0.99}$: for the first family,
the matching property gives $D_{vu}\le\Delta^{-0.99}W_v$.
Let $A_u$ indicate that some entry at $u$ activates.  These indicators
are independent over the hosts $u$, and
\begin{equation}\label{eq:host-activation-probability}
 \PP(A_u=1)\ge1-e^{-\tau W_u}\ge\tau K/4
\end{equation}
for large $\Delta$, since $K/2\le W_u\le2K$.
Lemma~\ref{lem:read-concentration} with $p=1$ and
$\sum_ub_u^2\le4K^2\Delta^{-0.98}$ gives
\begin{equation}\label{eq:activated-neighbor-mass}
 \sum_{u\in N_{G'}(v)}b_uA_u
 \ge\frac{\tau K}{4}\sum_{u\in N_{G'}(v)}b_u-\zeta
\end{equation}
except with probability at most $\exp(-\Delta^{0.5})$.
The next lemma bounds the loss caused by conflicting activations.

\begin{lemma}[Collision bound]\label{lem:activation-collisions}
Let $\mathcal E_v$ be the set of edges of $H$ with an endpoint hosted
at a neighbor of $v$ in $G'$, and let $X_v$ count the edges of
$\mathcal E_v$ whose two endpoints activate.  Then
\begin{equation}\label{eq:collision-tail}
 \PP\bigl(4K\Delta^{-0.99}X_v>\zeta\bigr)
 \le O(\Delta^3)e^{e-1-\Delta^{0.01}}+e^{-\Delta^{0.8}}.
\end{equation}
For either coefficient family in \eqref{eq:neighbor-coefficients},
the sum of $b_u$ over hosts $u$ counted by $A_u$ that receive no
color is at most $4K\Delta^{-0.99}X_v$.
\end{lemma}

\begin{proof}
A host with an activating entry but no successful entry contains an
endpoint of an edge counted by $X_v$.  Each such edge has at most
two hosts and each coefficient is at most $2K\Delta^{-0.99}$, proving
the final assertion.  Moreover,
\begin{equation}\label{eq:collision-expectation}
 \mu_v:=\EE X_v
 \le\tau^2\sum_{u\in N_{G'}(v)}
               \sum_{i\in L(u)\cap V(H)}w_is_i
 \le2K\Delta\tau^2\Lambda=\Delta^{0.80+o(1)}.
\end{equation}

For a possible endpoint $i$ of an edge in $\mathcal E_v$, the number
of its activating neighbors is a sum of independent indicators of
total mean $\tau s_i\le\tau\Lambda\le1$.  Its exponential moment at
parameter $1$ is at most $\exp(e-1)$, so it exceeds
$\Delta^{0.01}$ with probability at most
$\exp(e-1-\Delta^{0.01})$.  All possible endpoints are hosted within
distance two of $v$ in $G'$, and there are at most $k\le\Delta$ entries
at each host.  A union bound over $O(\Delta^3)$ endpoints shows that,
apart from the first error term in \eqref{eq:collision-tail}, the
graph formed by the edges counted by $X_v$ has maximum degree at
most $\Delta^{0.01}$.

If $4K\Delta^{-0.99}X_v>\zeta$, then
$X_v>\Delta^{0.84}/(4K)$.  In a graph of maximum degree at most
$\Delta^{0.01}$, greedy selection of disjoint edges produces a matching
of size at least $X_v/(2\Delta^{0.01})$.  For large $\Delta$ this is
at least $m=\lceil\Delta^{0.82}\rceil$.  For $e=il\in\mathcal E_v$
write $q_e=\tau^2w_iw_l$.  The events that both endpoints activate
are independent along any matching, because its edges have disjoint
endpoints and the activation flags are independent.  Consequently,
\begin{equation}\label{eq:collision-matching-tail}
 \begin{aligned}
 \PP\left(\begin{gathered}
     \text{some size-$m$ matching in $\mathcal E_v$}\\
     \text{has all endpoints activating}
   \end{gathered}\right)
 &\le\sum_{\substack{M\subseteq\mathcal E_v\\M\text{ matching},\ |M|=m}}
                      \prod_{e\in M}q_e\\
 &\le\frac{(\sum_{e\in\mathcal E_v}q_e)^m}{m!}
   =\frac{\mu_v^m}{m!}
   \le\left(\frac{e\mu_v}{m}\right)^m
   \le e^{-\Delta^{0.8}}.
 \end{aligned}
\end{equation}
The second inequality follows by expanding the power: each matching
appears $m!$ times, and all other terms are nonnegative.  The final
inequality uses \eqref{eq:collision-expectation} and
$m=\Delta^{0.82+o(1)}$; it is immediate if $\mu_v=0$.
Combining the two exceptional events proves the lemma.
\end{proof}

Together, \eqref{eq:activated-neighbor-mass} and
Lemma~\ref{lem:activation-collisions} imply that, for either choice
in \eqref{eq:neighbor-coefficients}, we may require
\begin{equation}\label{eq:colored-neighbor-mass}
 \sum_{\substack{u\in N_{G'}(v)\\u\text{ receives a color}}}b_u
 \ge\frac{\tau K}{4}\sum_{u\in N_{G'}(v)}b_u-2\zeta.
\end{equation}
Let $E_v$ be the event that \eqref{eq:vertex-concentration},
\eqref{eq:matching-concentration} for an edge incident with $v$, or
\eqref{eq:colored-neighbor-mass} for one of the two coefficient
families fails.  A union bound gives, for all sufficiently large
$\Delta$,
\begin{equation}\label{eq:round-failure-probability}
 \PP(E_v)\le\exp(-\Delta^{0.005}).
\end{equation}

\subsection{Simultaneous bounds and restoration of the weight conditions}

We use the Lov\'asz local lemma~\cite{ErdosLovasz1975} to avoid every
$E_v$.  Assign each multiplier flag to the host of a chosen center of
a radius-$R_0$ ball in $H$ containing its listed support.  Assign
activation and drop flags to the hosts of the activating and dropped
entries, respectively.  Every input determining $E_v$ is assigned
within distance $R_0+3$ of $v$ in $G'$.  Indeed, the host projection
does not increase distances; the provisional weights on an incident
matching use flags touching entries at its two hosts; and whether
a neighbor of $v$ receives a color is determined by activation flags
at distance at most two from $v$.  Events at distance greater than
$2(R_0+3)$ therefore involve disjoint sets of independent inputs.
A bound for the maximum degree of a dependency graph is
\begin{equation}\label{eq:round-dependency-degree}
 d=(\Delta+1)^{2R_0+10}
   \le\exp\bigl(O(\log\Delta\log\log\Delta)\bigr).
\end{equation}

For completeness, the symmetric local lemma says that finitely many
events of probability at most $q$, with a dependency graph of maximum
degree at most $d\ge1$, can all be avoided if
$q\le1/(e(d+1))$.  Put $x=1/(d+1)$, so
$q\le x(1-x)^d$.  Induction on the number of conditioned events shows
that any event, conditioned on the avoidance of any set of the others,
has probability at most $x$.  To make the induction step, split the
conditioning set into nonneighbors and at most $d$ neighbors of the
event.  Independence bounds the numerator after conditioning on the
nonneighbors by $q$, and the induction hypothesis bounds the
probability of avoiding the neighbors, under that conditioning, below
by $(1-x)^d$.  The same induction ensures that all conditioning
probabilities are positive.  Multiplying these positive conditional
avoidance probabilities proves joint avoidance.  By
\eqref{eq:round-failure-probability} and
\eqref{eq:round-dependency-degree}, its hypothesis holds here for
large $\Delta$.

\begin{proof}[Proof of Proposition~\ref{prop:one-round}]
Fix a trial avoiding all $E_v$.  Lemma~\ref{lem:round-expectation},
\eqref{eq:vertex-concentration}, and
\eqref{eq:matching-concentration} summed over the edges incident
with $v$ give
\begin{equation}\label{eq:provisional-bounds}
 \bar W_v\le W_v+o(\xi),\qquad
 \bar F_v\le F_v+\tau(B+1)D_v+o(\xi),
\end{equation}
because $\zeta+\Delta\Delta^{-1.25}=o(\xi)$.
Remove the vertices colored in this trial.  At a remaining vertex
$v$, only the edge part of its potential changes.  Each removed
neighbor $u$ contributes
$\bar D_{vu}\ge D_{vu}/2-\Delta^{-1.25}$ by
\eqref{eq:round-expectations} and
\eqref{eq:matching-concentration}.  Using
\eqref{eq:colored-neighbor-mass} with $b_u=D_{vu}$, the total
decrease is at least
\begin{equation}\label{eq:potential-removal-gain}
 \sum_{\substack{u\in N_{G'}(v)\\u\text{ receives a color}}}\bar D_{vu}
 \ge\frac{\tau K}{8}D_v-\zeta-\Delta^{-0.25}.
\end{equation}
Since $K/8>B+1$, the potential after removal is at most $F_v+o(\xi)$.
The same requirement with $b_u=1/\Delta$ gives the degree bound in
\eqref{eq:round-conclusions}, because $K/4\ge1$ and $2\zeta=o(\xi)$.
It remains to restore \eqref{eq:weight-bounds} without losing the
mass and potential estimates.

To \emph{trim} an entry means to set its weight to zero.  If its
current positive weight is $z$, the change in the non-edge part of
its host's potential is exactly
\begin{equation}\label{eq:trimming-cost}
 c(z):=z\bigl(1-\log(z/w_*)+\gamma\log\Delta\bigr).
\end{equation}
First trim all entries with $z>\Delta^{-0.99}$.  Their costs are
nonpositive for large $\Delta$, since
$\log(\Delta^{-0.99}/w_*)=0.01\log\Delta-O(\log\log\Delta)$ and
$\gamma<0.01$.  Next trim all positive entries with $z<\Delta^{-3}$.
The function $c(z)$ is increasing on this interval for large $\Delta$
and has limit zero at zero, so each cost is at most
$O(\Delta^{-3}\log\Delta)$.  There are at most $k\le\Delta$ entries
at a host, giving total cost $o(\xi)$ there.  Both trimming operations
can only decrease edge terms.

At the resulting weights, mark each positive entry $i$ with
$s_i>\Lambda$, and trim all marked entries simultaneously.  For a
marked entry of weight $z$, \eqref{eq:trimming-cost} is at most
$\Lambda z$ on the restored weight interval.  To account for the
edge loss without double counting, let $M$ be the marked entries and
fix a host $v$.  In the next identity, $H,w_i,s_i,D_v$ refer to the
remaining graph and weights immediately before this last trim, and
$D'_v$ is the edge term after it.  Then
\begin{equation}\label{eq:simultaneous-trimming-loss}
 D_v-D'_v
 =\sum_{i\in M\cap L(v)}w_is_i
  +\sum_{i\in (V(H)\setminus M)\cap L(v)}
       w_i\sum_{l\in M\cap N_H(i)}w_l
 \ge\sum_{i\in M\cap L(v)}w_is_i.
\end{equation}
Each deleted edge in $D_v$ has exactly one endpoint hosted at $v$, so
it occurs in precisely one of these two sums.  The first sum pays for
all non-edge costs at $v$, since $w_is_i>\Lambda w_i$ for marked
entries.  This remains true when both ends of an edge are trimmed.
This last trim does not
increase any $F_v$.  The unmarked entries had $s_i\le\Lambda$ before
the trim, and deleting weights only reduces their neighborhood
weights, so the remaining entries satisfy \eqref{eq:weight-bounds}.

All trims preserve availability and do not increase $W_v$.  Combining
their costs with \eqref{eq:provisional-bounds} and
\eqref{eq:potential-removal-gain} proves the mass and potential bounds
in \eqref{eq:round-conclusions} for sufficiently large $\Delta$.
\end{proof}
